\documentclass[12pt,a4paper]{amsart}

\usepackage{amsmath,amssymb,amsthm}
\usepackage{mathtools}
\usepackage{hyperref}
\usepackage{geometry}
\geometry{margin=2.5cm}

% -------------------------------------------------------
% Theorem environments
% -------------------------------------------------------
\newtheorem{theorem}{Theorem}[section]
\newtheorem{lemma}[theorem]{Lemma}
\newtheorem{corollary}[theorem]{Corollary}
\newtheorem{proposition}[theorem]{Proposition}
\theoremstyle{definition}
\newtheorem{definition}[theorem]{Definition}
\newtheorem{remark}[theorem]{Remark}

% -------------------------------------------------------
% Custom commands
% -------------------------------------------------------
\newcommand{\N}{\mathbb{N}}
\newcommand{\Z}{\mathbb{Z}}
\newcommand{\Q}{\mathbb{Q}}
\newcommand{\Fp}{\mathbb{F}_p}
\DeclareMathOperator{\ord}{ord}
\DeclareMathOperator{\lcm}{lcm}

% -------------------------------------------------------
\begin{document}
% -------------------------------------------------------

\title{Wilson's Theorem for Prime Powers}

\author{Michael Fuhrmann}
\address{Specialist Mathematics Project, Build 63}
\email{specialist-maths@localhost}
\date{\today}

\begin{abstract}
We investigate the product of all units in the ring $\Z/p^k\Z$ and prove the
following complete result: for an odd prime $p$ and any $k \geq 1$, the product
of all integers in $\{1, \ldots, p^k\}$ coprime to $p$ satisfies
\[
  \prod_{\substack{a=1\\\gcd(a,p^k)=1}}^{p^k} a \;\equiv\; -1 \pmod{p^k}.
\]
For $p = 2$ the situation is exceptional: the analogous product equals $1 \pmod{2^k}$
for $k \geq 3$, while the small cases $k = 1$ and $k = 2$ are handled separately.
The proof rests on two pillars: (i) the pairing method — every unit pairs with its
inverse — and (ii) the fact that the only square roots of unity modulo $p^k$
(for odd $p$) are $\pm 1$.  As a corollary we identify the structure of the
unit group $(\Z/p^k\Z)^*$.
\end{abstract}

\maketitle
\tableofcontents

% ================================================================
\section{Introduction}
% ================================================================

Wilson's classical theorem states that $(p-1)! \equiv -1 \pmod{p}$ for every
prime $p$.  The product $(p-1)!$ is nothing other than the product of \emph{all}
elements of the multiplicative group $(\Z/p\Z)^*$.
A natural question is: what happens when we replace the modulus $p$ by a
prime power $p^k$?

The unit group $(\Z/p^k\Z)^*$ consists of all residues $1 \leq a \leq p^k$
with $\gcd(a, p^k) = 1$, and has order $\phi(p^k) = p^{k-1}(p-1)$.
For odd primes $p$ this group is cyclic, and the product of all its elements
can be computed from the group structure alone.
For $p = 2$ the group is no longer cyclic for $k \geq 3$, and the product
turns out to be $+1$ rather than $-1$.

Throughout, $p$ denotes a prime and $k$ a positive integer.  All congruences
are modulo $p^k$ unless stated otherwise.

% ================================================================
\section{Square Roots of Unity Modulo \texorpdfstring{$p^k$}{p\^{}k}}
% ================================================================

The key lemma identifies the self-inverse elements of $(\Z/p^k\Z)^*$.

\begin{lemma}[Square roots of unity for odd prime powers]%
  \label{lem:sqrt_unity}
  Let $p$ be an odd prime and $k \geq 1$.  The solutions of
  $a^2 \equiv 1 \pmod{p^k}$ in $\{1, \ldots, p^k\}$ are exactly
  $a = 1$ and $a = p^k - 1 \equiv -1 \pmod{p^k}$.
\end{lemma}

\begin{proof}
  We have $a^2 \equiv 1 \pmod{p^k}$ if and only if $p^k \mid (a-1)(a+1)$.
  Since $\gcd(a-1, a+1) \mid 2$ and $p$ is odd, $p$ cannot divide both
  $a-1$ and $a+1$ simultaneously.  Therefore $p^k \mid a-1$ or $p^k \mid a+1$,
  which gives $a \equiv 1 \pmod{p^k}$ or $a \equiv -1 \pmod{p^k}$.

  More precisely: suppose $p^k \mid (a-1)(a+1)$.  Write $a - 1 = p^\alpha u$
  and $a + 1 = p^\beta v$ with $\gcd(u,p) = \gcd(v,p) = 1$.
  Then $(a+1) - (a-1) = 2$, so $p^{\min(\alpha,\beta)} \mid 2$.
  Since $p$ is odd, $\min(\alpha, \beta) = 0$.
  Hence either $\alpha = 0$ (meaning $p \nmid a-1$) and $p^k \mid a+1$,
  or $\beta = 0$ (meaning $p \nmid a+1$) and $p^k \mid a-1$.
  Both cases yield exactly the solutions $a \equiv \pm 1 \pmod{p^k}$.
\end{proof}

\begin{remark}
  Lemma~\ref{lem:sqrt_unity} fails for $p = 2$ and $k \geq 3$:
  the congruence $a^2 \equiv 1 \pmod{8}$ has four solutions modulo $8$:
  $a \in \{1, 3, 5, 7\}$.  In general, for $k \geq 3$, $a^2 \equiv 1
  \pmod{2^k}$ has four solutions: $a \equiv \pm 1 \pmod{2^{k-1}}$.
\end{remark}

% ================================================================
\section{The Main Theorem for Odd Primes}
% ================================================================

\begin{theorem}[Wilson for prime powers, odd $p$]\label{thm:wilson_pk_odd}
  Let $p$ be an odd prime and $k \geq 1$.  Then
  \[
    \prod_{\substack{a=1\\\gcd(a,p^k)=1}}^{p^k} a
    \;\equiv\; -1 \pmod{p^k}.
  \]
\end{theorem}

\begin{proof}
  Denote by $P$ the product on the left-hand side.  We work in the group
  $G = (\Z/p^k\Z)^*$ of order $\phi(p^k) = p^{k-1}(p-1)$.

  \textbf{Step 1 (Pairing).}
  For every $a \in G$, pair $a$ with its multiplicative inverse
  $a^{-1} \pmod{p^k}$.  If $a \neq a^{-1}$, then $\{a, a^{-1}\}$ is a
  two-element subset of $G$ and contributes the factor $a \cdot a^{-1} = 1$
  to $P$.

  \textbf{Step 2 (Self-inverse elements).}
  An element $a$ is self-inverse (i.e.\ $a = a^{-1}$) if and only if
  $a^2 \equiv 1 \pmod{p^k}$.  By Lemma~\ref{lem:sqrt_unity}, the only
  self-inverse elements are $a \equiv 1$ and $a \equiv -1 \pmod{p^k}$.

  \textbf{Step 3 (Computing the product).}
  After pairing, all elements of $G \setminus \{1, -1\}$ pair up into
  $(|G| - 2)/2 = (p^{k-1}(p-1) - 2)/2$ pairs, each contributing factor $1$.
  The remaining factors are those from the two self-inverse elements:
  \[
    P \;\equiv\; 1 \cdot (-1) \;\equiv\; -1 \pmod{p^k}.
  \]

  \textbf{Step 4 (Validity of the pairing count).}
  We must verify that $|G| - 2$ is even, i.e.\ that $|G| = p^{k-1}(p-1)$
  is even.  Since $p$ is an odd prime, $p - 1$ is even, so $p^{k-1}(p-1)$
  is indeed even.  Therefore the pairing is valid and the calculation above
  is correct.
\end{proof}

% ================================================================
\section{The Case \texorpdfstring{$p = 2$}{p = 2}}
% ================================================================

The prime $p = 2$ requires a separate analysis because the group
$(\Z/2^k\Z)^*$ is not cyclic for $k \geq 3$.

\begin{theorem}[Wilson for powers of $2$]\label{thm:wilson_pk_2}
  Let $k \geq 1$.  Then
  \[
    \prod_{\substack{a=1\\\gcd(a,2^k)=1}}^{2^k} a
    \;\equiv\;
    \begin{cases}
      1  \pmod{2}   & k = 1, \\
      -1 \pmod{4}   & k = 2, \\
      1  \pmod{2^k} & k \geq 3.
    \end{cases}
  \]
\end{theorem}

\begin{proof}
  \textbf{Case $k = 1$.}
  The only odd residue modulo $2$ in $\{1\}$ is $1$.  The product is $1$,
  and $1 \equiv 1 \pmod 2$.

  \textbf{Case $k = 2$ (modulo 4).}
  The units modulo $4$ are $\{1, 3\}$.  Their product is $1 \cdot 3 = 3 \equiv -1 \pmod 4$.

  \textbf{Case $k \geq 3$.}
  The unit group $(\Z/2^k\Z)^*$ has order $\phi(2^k) = 2^{k-1}$.
  For $k \geq 3$, it is \emph{not} cyclic; instead
  \[
    (\Z/2^k\Z)^* \;\cong\; \Z/2 \times \Z/2^{k-2}.
  \]
  The elements of order dividing $2$ (i.e.\ solutions to $a^2 \equiv 1 \pmod{2^k}$)
  are $a \equiv \pm 1 \pmod{2^{k-1}}$, giving four solutions modulo $2^k$:
  $1$, $2^{k-1}-1$, $2^{k-1}+1$, and $2^k - 1 \equiv -1$.

  Hence the subgroup of self-inverse elements is
  $S = \{1,\; -1,\; 2^{k-1}-1,\; 2^{k-1}+1\} \pmod{2^k}$,
  which has order $4$.  The product of all elements of $S$ is:
  \[
    1 \cdot (-1) \cdot (2^{k-1}-1) \cdot (2^{k-1}+1)
    = (-1)\bigl((2^{k-1})^2 - 1\bigr)
    = -(2^{2k-2} - 1).
  \]
  We compute this modulo $2^k$: since $k \geq 3$, we have $2k-2 \geq k+1 > k$,
  so $2^{2k-2} \equiv 0 \pmod{2^k}$, and
  \[
    -(2^{2k-2}-1) \equiv -(-1) = 1 \pmod{2^k}.
  \]
  All non-self-inverse elements pair into pairs $\{a, a^{-1}\}$ each
  contributing factor $1$.  Therefore the total product $P \equiv 1 \pmod{2^k}$.
\end{proof}

% ================================================================
\section{Structure of the Unit Group and Corollaries}
% ================================================================

The proof of Theorem~\ref{thm:wilson_pk_odd} relies on the cyclicity of
$(\Z/p^k\Z)^*$ for odd $p$, which we now state precisely.

\begin{theorem}[Structure of $(\Z/p^k\Z)^*$]\label{thm:structure}
  \begin{enumerate}
    \item[(a)] For every odd prime $p$ and $k \geq 1$:
      $(\Z/p^k\Z)^* \cong \Z_{\phi(p^k)} = \Z_{p^{k-1}(p-1)}$
      (cyclic of order $p^{k-1}(p-1)$).
    \item[(b)] For $p = 2$: $(\Z/2\Z)^* = \{1\}$, $(\Z/4\Z)^* \cong \Z_2$,
      and $(\Z/2^k\Z)^* \cong \Z_2 \times \Z_{2^{k-2}}$ for $k \geq 3$.
  \end{enumerate}
\end{theorem}

\begin{proof}
  \textbf{Part (a).}
  We show that $(\Z/p^k\Z)^*$ is cyclic by lifting a primitive root modulo $p$
  to a primitive root modulo $p^k$.

  Let $g$ be a primitive root modulo $p$ (which exists since $(\Z/p\Z)^*$ is
  cyclic of order $p-1$).  We may choose $g$ such that $g^{p-1} \not\equiv 1
  \pmod{p^2}$ (if this fails for $g$, replace $g$ by $g + p$, which still
  generates $(\Z/p\Z)^*$; at most one value fails since there are $p-1$ choices).

  We claim that $g$ is a primitive root modulo $p^k$ for all $k \geq 1$.
  By induction: suppose $g$ has order $p^{k-1}(p-1)$ modulo $p^k$.
  Then $g^{p^{k-1}(p-1)} \equiv 1 \pmod{p^k}$.  Write
  \[
    g^{p-1} = 1 + p \cdot t, \quad t \not\equiv 0 \pmod p.
  \]
  By the binomial theorem (or Hensel lifting),
  $g^{p^{k-1}(p-1)} = (1 + pt)^{p^{k-1}} \equiv 1 + p^k t' \pmod{p^{k+1}}$
  with $t' \not\equiv 0 \pmod p$.
  Hence the order of $g$ modulo $p^{k+1}$ is $p^k(p-1) = \phi(p^{k+1})$,
  proving $g$ is also a primitive root modulo $p^{k+1}$.

  \textbf{Part (b).}
  The unit group $(\Z/2\Z)^* = \{1\}$ is trivial.
  $(\Z/4\Z)^* = \{1, 3\} \cong \Z_2$.
  For $k \geq 3$: the element $-1 \equiv 2^k - 1$ generates a subgroup of
  order $2$, and $5$ generates a subgroup of order $2^{k-2}$ (since
  $5^{2^{k-3}} \equiv 1 + 2^{k-1} \not\equiv 1 \pmod{2^k}$ but
  $5^{2^{k-2}} \equiv 1 \pmod{2^k}$).
  Together $-1$ and $5$ generate the full group, and
  $\langle -1 \rangle \cap \langle 5 \rangle = \{1\}$, yielding the direct
  product decomposition $(\Z/2^k\Z)^* \cong \Z_2 \times \Z_{2^{k-2}}$.
\end{proof}

\begin{corollary}[Unique element of order 2 for odd prime powers]%
  \label{cor:unique_order2}
  For an odd prime $p$ and $k \geq 1$, the group $(\Z/p^k\Z)^*$ has exactly
  one element of order $2$, namely $-1 \equiv p^k - 1 \pmod{p^k}$.
\end{corollary}

\begin{proof}
  By Theorem~\ref{thm:structure}(a), $(\Z/p^k\Z)^*$ is cyclic of even order
  $p^{k-1}(p-1)$.  In any cyclic group of even order, there is exactly one
  element of order $2$, namely the unique element $\tau$ such that $\tau^2 = e$
  and $\tau \neq e$.  By Lemma~\ref{lem:sqrt_unity}, this element is $-1$.
\end{proof}

\begin{corollary}[Wilson product formula]\label{cor:wilson_formula}
  For an odd prime $p$ and $k \geq 1$:
  \[
    \prod_{\substack{a=1\\\gcd(a,p^k)=1}}^{p^k} a \;\equiv\; -1 \pmod{p^k}.
  \]
  For $k = 1$ this specialises to the classical Wilson theorem $(p-1)! \equiv -1 \pmod p$.
\end{corollary}

\begin{proof}
  This is Theorem~\ref{thm:wilson_pk_odd}.  For $k = 1$, the product over
  $\{a : 1 \leq a \leq p,\ \gcd(a,p) = 1\}$ equals $(p-1)!$.
\end{proof}

% ================================================================
\section{Conclusion}
% ================================================================

We have proved a complete generalisation of Wilson's theorem to prime powers.
The key insight is that for odd prime powers $p^k$, the group of units is
cyclic and has \emph{exactly one} self-inverse element other than the identity,
namely $-1$.  This forces the product of all units to be $-1$.

For $p = 2$, the breakdown of cyclicity for $k \geq 3$ leads to four
self-inverse elements whose collective product is $+1$, reversing the sign
compared to the odd-prime case.

\begin{thebibliography}{9}

\bibitem{Hardy1979}
G.~H.~Hardy and E.~M.~Wright.
\newblock \emph{An Introduction to the Theory of Numbers}, 5th~ed.
\newblock Oxford University Press, 1979.

\bibitem{Ireland1990}
K.~Ireland and M.~Rosen.
\newblock \emph{A Classical Introduction to Modern Number Theory}, 2nd~ed.
\newblock Springer, New York, 1990.

\bibitem{Niven1991}
I.~Niven, H.~S.~Zuckerman, and H.~L.~Montgomery.
\newblock \emph{An Introduction to the Theory of Numbers}, 5th~ed.
\newblock Wiley, New York, 1991.

\bibitem{Rose1988}
H.~E.~Rose.
\newblock \emph{A Course in Number Theory}, 2nd~ed.
\newblock Oxford University Press, 1994.

\end{thebibliography}

\end{document}
